% The rule-of-thumb controller as formulas, and its attachment to the first principles model
% inside the augmentation training loop.
% Notation follows docs/writeup/jan-augmentation-writeup.tex and the sibling note
% scripts/gantry/closed-loop-controller/controller-in-derivation.tex.
% Compile: pdflatex controller-in-training.tex
\documentclass[11pt]{article}

\usepackage[a4paper,margin=2.2cm]{geometry}
\usepackage[T1]{fontenc}
\usepackage[utf8]{inputenc}
\usepackage{lmodern}
\usepackage{amsmath,amssymb,mathtools,bm}
\usepackage{microtype}
\usepackage{enumitem}
\usepackage{titlesec}
\usepackage{booktabs}
\usepackage{graphicx}
\usepackage[hidelinks]{hyperref}

\titleformat{\section}{\Large\bfseries}{\thesection}{0.75em}{}
\titleformat{\subsection}{\large\bfseries}{\thesubsection}{0.75em}{}
\titlespacing*{\section}{0pt}{1.4ex plus .2ex minus .2ex}{0.9ex}
\titlespacing*{\subsection}{0pt}{1.0ex plus .2ex minus .2ex}{0.5ex}

\newcommand{\R}{\mathbb{R}}
\newcommand{\diag}{\operatorname{diag}}
\newcommand{\Cfb}{C_{\mathrm{fb}}}
\newcommand{\Cnorm}{C_{\mathrm{n}}}
\newcommand{\wb}{\omega_{\mathrm{b}}}
\newcommand{\fms}{f_{\mathrm{ms}}}
\newcommand{\ufb}{u_{\mathrm{fb}}}
\newcommand{\udata}{u_{\mathrm{data}}}
\newcommand{\umodel}{u_{\mathrm{model}}}
\newcommand{\ydata}{y_{\mathrm{data}}}
\newcommand{\ymodel}{y_{\mathrm{model}}}
\newcommand{\Yop}{Y_{\mathrm{op}}}
\newcommand{\Tu}{T_{u}}
\newcommand{\Ty}{T_{y}}

\setlist{nosep, leftmargin=1.4em}

\begin{document}

\begin{center}
  {\LARGE\bfseries The feedback controller in the training loop}\\[0.4em]
  {\large Dirk van den Berg}\\[0.2em]
  {\small Section~\ref{sec:formulas} writes the MATLAB rule-of-thumb controller as formulas.
  Section~\ref{sec:attach} attaches it to the first principles model. Section~\ref{sec:block}
  states that it is a separate training block.}
\end{center}
\vspace{0.6em}

\noindent\textbf{Symbols.}\quad The baseline write-up already uses $C$ for the damping matrix and
$K$ for the stiffness matrix, and the sibling note \texttt{controller-in-derivation.tex} uses
$S_{\mathrm{o}}$ for the output sensitivity. The controller, its design gain and the normalisation
matrices therefore take different symbols here.

\begin{center}
\begin{tabular}{lll}
\toprule
symbol & meaning & would clash with \\
\midrule
$\Cfb(z)$        & discrete feedback controller, $3\times3$ diagonal & $C$, damping matrix \\
$\Cnorm(s)$      & the shared normalised controller shape            & \\
$\kappa_j$       & per-channel normalisation gain                    & $K$, stiffness matrix \\
$\wb$            & design bandwidth in rad/s                         & \\
$C_{\mathrm{d}}$ & output matrix of the discrete model               & $C$, damping matrix \\
$\Tu,\ \Ty$      & input and output normalisation, diagonal          & $S_{\mathrm{o}}$, sensitivity \\
\bottomrule
\end{tabular}
\end{center}

\section{The rule-of-thumb controller as formulas}
\label{sec:formulas}

The controller is not our design. It is \texttt{ruleOfThumb.m} of the read-only first principles
model, called once per channel from \texttt{gtd\_build\_plant.m}. This section writes each line of
that function as a formula.

\subsection{The design plant}

The controller is designed on a plant that is stage in and stage out, frozen at the record's
operating point $\Yop$ and carrying the full payload mass, so that it is linear time invariant
even though the true plant is not:
\begin{equation}
  \mathrm{sys}(\Yop) \;=\; P^{\top}\,
  G\bigl(M_{\mathrm{op}}(\Yop),\, C_{\mathrm{damp}},\, K\bigr)\, P ,
  \qquad
  P =
  \begin{bmatrix}
    1 & 1 & 0 \\[1mm]
    \tfrac{L_b}{2} & -\tfrac{L_b}{2} & 0 \\[1mm]
    0 & 0 & 1
  \end{bmatrix} .
  \label{eq:sys}
\end{equation}
The loop is therefore diagonal in the \emph{stage} frame $[\,X_1, X_2, Y\,]$, and the logical
frame appears only inside $\mathrm{sys}$. Only the three diagonal channels
$\mathrm{sys}_{jj}(s)$, $j = 1,2,3$, are ever used.

\subsection{The shape}

Let $f_{\mathrm{bw}} = 100$~Hz and $\wb = 2\pi f_{\mathrm{bw}}$. All three channels share one
shape, built from the three factors of \texttt{ruleOfThumb.m}:
\begin{equation}
  \underbrace{\frac{s+\wb/6}{s}}_{\texttt{int}}
  \qquad
  \underbrace{\frac{s+\wb/3}{s+3\wb}}_{\texttt{leadlag}}
  \qquad
  \underbrace{\frac{10\,\wb}{s+10\,\wb}}_{\texttt{lowpass}}
  \label{eq:factors}
\end{equation}
\begin{equation}
  \Cnorm(s)
  \;=\;
  \frac{10\,\wb\,\bigl(s+\tfrac{1}{6}\wb\bigr)\bigl(s+\tfrac{1}{3}\wb\bigr)}
       {s\,\bigl(s+3\wb\bigr)\bigl(s+10\,\wb\bigr)} .
  \label{eq:cnorm}
\end{equation}

\subsection{The gain}

The plant enters the design through one scalar per channel only. The gain is fixed by requiring
unit open-loop magnitude at the bandwidth,
\begin{equation}
  \kappa_j \;=\; \frac{1}{\bigl|\,\mathrm{sys}_{jj}(i\wb)\,\Cnorm(i\wb)\,\bigr|} ,
  \qquad j = 1,2,3 ,
  \label{eq:kappa}
\end{equation}
so that $f_{\mathrm{bw}}$ is the $0$~dB crossover of each loop by construction. Equation
\eqref{eq:kappa} is the entire $\Yop$ dependence of the controller: three numbers.

\subsection{The discretisation}

\begin{equation}
  C_j(z) \;=\; \kappa_j\,\Cnorm(s)\Bigr|_{\,s \,=\, \frac{2}{T_s}\frac{z-1}{z+1}} ,
  \qquad
  \Cfb(z) \;=\; \diag\bigl(C_1(z),\, C_2(z),\, C_3(z)\bigr) .
  \label{eq:cfb}
\end{equation}

\subsection{The realisation used in the loop}

Each channel is of order three, so the block diagonal realisation has nine states. It is
biproper:
\begin{equation}
  \Cfb : \quad
  \begin{aligned}
    x_{c,k+1} &= A_c\,x_{c,k} + B_c\,e_{c,k} , \\
    \ufb{}_{,k}   &= C_c\,x_{c,k} + D_c\,e_{c,k} ,
  \end{aligned}
  \qquad
  x_c \in \R^{9} , \qquad D_c \neq 0 .
  \label{eq:cfbss}
\end{equation}

\section{Attaching the controller to the first principles model}
\label{sec:attach}

\subsection{The model}

The augmented baseline in discrete time, normalised coordinates, with the first principles block
on the physical rows and the learned block added on the routed rows:
\begin{equation}
  x_{k+1} = f(x_k, u_k) ,
  \qquad
  y_k = h(x_k) = C_{\mathrm{d}}\,x_k + b ,
  \qquad
  x \in \R^{8},\; u \in \R^{3},\; y \in \R^{3} .
  \label{eq:model}
\end{equation}
The output map is affine in the state and has no feedthrough,
\begin{equation}
  \frac{\partial h}{\partial u} \;=\; 0
  \qquad (D_{\mathrm{d}} = 0) .
  \label{eq:Dd}
\end{equation}

\subsection{The residual form}

The record was generated by the machine loop, and the same controller is closed around the model
from the same reference and the same injected multisine:
\begin{equation}
  \udata = \Cfb\bigl(r - \ydata\bigr) + \fms ,
  \qquad
  \umodel = \Cfb\bigl(r - \ymodel\bigr) + \fms .
  \label{eq:twoloops}
\end{equation}
Subtracting the two, with $r$ and $\fms$ common to both and $\Cfb$ linear,
\begin{equation}
  \boxed{\;
    \umodel{}_{,k}
    \;=\;
    \udata{}_{,k}
    \;+\;
    \Cfb\bigl(\ydata - \ymodel\bigr)_k \; }
  \label{eq:residual}
\end{equation}
Neither $r$ nor $\fms$ survives the subtraction, so the training loop needs only the recorded
input and the recorded output.

\subsection{The stepped recursion}

With $\Tu = \diag(u_{\mathrm{std}})$ and $\Ty = \diag(y_{\mathrm{std}})$ the normalisation
matrices, equation \eqref{eq:residual} is realised as
\begin{equation}
  \begin{aligned}
    \ymodel{}_{,k} &= h(x_k) , \\
    e_{c,k}        &= \Ty\bigl(\ydata{}_{,k} - \ymodel{}_{,k}\bigr) , \\
    \ufb{}_{,k}    &= C_c\,x_{c,k} + D_c\,e_{c,k} , \\
    u_k            &= \udata{}_{,k} + \Tu^{-1}\,\ufb{}_{,k} , \\
    x_{k+1}        &= f(x_k,\, u_k) , \\
    x_{c,k+1}      &= A_c\,x_{c,k} + B_c\,e_{c,k} ,
  \end{aligned}
  \qquad
  \begin{aligned}
    &x_{\tau|\tau}   = \mathrm{encoder}\bigl(u_{\tau-n_b:\tau+n_{b,r}},\;
                                             y_{\tau-n_a:\tau+n_{a,r}}\bigr) , \\[1mm]
    &x_{c,\tau|\tau} = 0 .
  \end{aligned}
  \label{eq:recursion}
\end{equation}
The error $e_{c}$ is in metres and $\ufb$ is in newtons, so the two normalisation matrices convert
into and out of the physical units the controller was designed in. The input mean does not appear
because $\Tu^{-1}\ufb$ is an increment on the recorded input.

\medskip\noindent\textbf{The ordering of \eqref{eq:recursion} is forced.}
\begin{itemize}
  \item $D_c \neq 0$ in \eqref{eq:cfbss}: the controller output at step $k$ depends on the error
        at step $k$, so $\ymodel{}_{,k}$ has to be computed before $u_k$.
  \item $D_{\mathrm{d}} = 0$ in \eqref{eq:Dd}: the model output at step $k$ does not depend on
        $u_k$, so the loop closes without an algebraic loop, at one model evaluation per step.
\end{itemize}

\section{The controller is a separate training block}
\label{sec:block}

Equation \eqref{eq:recursion} steps two subsystems in sequence rather than joining them into one
state vector. The alternative, joining the controller into the interconnect, would carry the
augmented state $z = [\,x^{\top}\; x_c^{\top}\,]^{\top} \in \R^{17}$ and the single map
\begin{equation}
  z_{k+1} =
  \begin{bmatrix}
    f\Bigl(x_k,\;\udata{}_{,k}
      + \Tu^{-1}\bigl(C_c x_{c,k} + D_c\,\Ty(\ydata{}_{,k} - h(x_k))\bigr)\Bigr) \\[2mm]
    A_c\,x_{c,k} + B_c\,\Ty\bigl(\ydata{}_{,k} - h(x_k)\bigr)
  \end{bmatrix} .
  \label{eq:joined}
\end{equation}
The separate form was chosen. In it the trainable parameters enter only through the model,
\begin{equation}
  \theta \;\text{ appears only in }\; f \;\text{ and }\; h ,
  \qquad
  (A_c, B_c, C_c, D_c) \;\text{ constant} ,
  \label{eq:theta}
\end{equation}
with four consequences:
\begin{enumerate}
  \item The model state stays at $n_x = 8$ rather than $17$, so the encoder keeps its width and
        its meaning.
  \item The encoder never has to reconstruct a controller state. By \eqref{eq:recursion},
        $x_{c,\tau|\tau} = 0$, so there is nothing to encode.
  \item The Jacobian of the first principles block with respect to $\theta$ is not premultiplied
        by the normalisation and controller matrices of \eqref{eq:joined}. Gradients still pass
        through $\Cfb$, but only as a linear filter with fixed coefficients.
  \item $\Cfb$ may differ per record. It is rebuilt at each record's $\Yop$ through
        \eqref{eq:sys} and \eqref{eq:kappa}, giving nine distinct instances across the
        22 records, selected inside the batch by a per-sample record index. The joined form
        \eqref{eq:joined} would have to be rebuilt for every record.
\end{enumerate}

\subsection{The one approximation}

Everything above is an identity except the initial condition $x_{c,\tau|\tau} = 0$. In the
residual form the controller filters $\ydata - \ymodel$, which does not exist before the
truncated window opens at $\tau$, because the model trajectory starts at $\tau$. The value is
therefore a definition rather than an estimate of an unknown, and zero is the unique choice for
which the correction term in \eqref{eq:residual} vanishes identically when the model is exact.

What it costs is integral memory. $\Cnorm$ has a pole at $s = 0$, so $\Cfb$ has a pole at $z = 1$
and never forgets DC: run continuously, $x_c$ converges to the constant force that nulls the DC
part of the model error, while resetting it at every window start discards that. The validation
free run does accumulate it over the whole record, so the reset is a training against validation
asymmetry, and it is measured rather than argued.

\section{Verification}

[figures]

\end{document}
